\section{Two dispersion thresholds and the selection of critical profiles}
\label{sec:r10-symbols}

The loss of a profile gap under a small nonhomogeneous perturbation has a
more precise variational description.  We prove two distinct mass
thresholds, classify the upper endpoint for a full interpolation family,
and determine the scale and selected profiles as the lower threshold is
approached.  The selection argument requires neither uniqueness of the
homogeneous optimizers nor nondegeneracy of their real linearization.
These are constrained variational statements; they do not classify
arbitrary solutions of the homogeneous evolution.

\subsection{Symbols and the two thresholds}

Fix $N\ge1$, $m>1$, $s=m/2$, $r=m/N$, and $p=2+2r$.  Let $a_0$ and
$a_\infty$ be real, even, positive elliptic symbols homogeneous of degree
$m$, continuous on the sphere and $C^{N+2}$ off zero.  Suppose
$a_\infty-a_0$ is also positive elliptic.  Use the notation
\[
 M(f)=\|f\|_2^2,\qquad A_j(f)=\int a_j(\xi)|\widehat f(\xi)|^2\,d\xi,
 \qquad E_j(f)=\tfrac12 A_j(f)-\tfrac1p\|f\|_p^p
 \quad(j=0,\infty).
\]
The sharp homogeneous masses $M_j$ are defined by
\begin{equation}\label{eq:r10-sharp-gn}
 \frac1p\|f\|_p^p\le
 \frac12\left(\frac{M(f)}{M_j}\right)^r A_j(f).
\end{equation}
In particular, $M_\infty>M_0$: if
$a_\infty\ge(1+\eta)a_0$, sharpness gives
$M_\infty\ge(1+\eta)^{1/r}M_0$.

Write the full symbol as $k=a_0+b$, with $b$ real, even, and continuous
off zero.  Assume that, for some $q>0$ and constants $c,C>0$,
\begin{equation}\label{eq:r10-bounds}
 c\min\{|\xi|^{m+q},|\xi|^m\}
 \le b(\xi)\le C\min\{|\xi|^{m+q},|\xi|^m\},
\end{equation}
and that, uniformly for $\theta\in S^{N-1}$,
\begin{equation}\label{eq:r10-high-limit}
 R^{-m}k(R\theta)\longrightarrow a_\infty(\theta)
 \qquad(R\longrightarrow\infty).
\end{equation}
Thus $k$ is comparable with $|\xi|^m$.  Define
\[
 B(f)=\int b(\xi)|\widehat f(\xi)|^2\,d\xi,\qquad
 E_k(f)=\tfrac12[A_0(f)+B(f)]-\tfrac1p\|f\|_p^p,
 \qquad I(M)=\inf_{\substack{f\in H^s\\M(f)=M}}E_k(f).
\]

\begin{theorem}[Two homogeneous limits determine two mass thresholds]
\label{thm:r10-two-thresholds}
Under \eqref{eq:r10-bounds}--\eqref{eq:r10-high-limit}, the sharp GN
mass for the full form $A_0+B$ is exactly $M_0$, and its sharp inequality
has no nonzero optimizer.  The constrained infima satisfy
\begin{enumerate}
 \item $I(M)=0$, without attainment, for $0<M\le M_0$;
 \item $-\infty<I(M)<0$ for $M_0<M<M_\infty$. Every minimizing
 sequence has, after translations and extraction, a strong $H^s$ limit
 attaining $I(M)$. A real minimizer can be chosen;
 \item $I(M)=-\infty$ for $M>M_\infty$.
\end{enumerate}
The symbols $a_0$ and $a_\infty$ need not have the same anisotropy.
\end{theorem}

\begin{proof}
Since $A_0+B\ge A_0$, the full GN constant is at most the principal
one. If $Q$ is a principal optimizer and
$Q_\lambda(x)=\lambda^{N/2}Q(\lambda x)$, then
\[
 \lambda^{-m}B(Q_\lambda)
 =\int\lambda^{-m}b(\lambda\xi)|\widehat Q(\xi)|^2\,d\xi
 \longrightarrow0\qquad(\lambda\downarrow0)
\]
by dominated convergence and \eqref{eq:r10-bounds}. Thus the sharp
constants coincide. Since $B(f)>0$ for every nonzero $f$, equality in
the full inequality is impossible. At mass $M\le M_0$, the sharp
principal inequality gives $E_k(f)>0$. Dilating any smooth mass-$M$
trial to low frequency proves that the infimum is zero.

For every $\eta>0$, \eqref{eq:r10-high-limit} gives a finite $C_\eta$
such that $k\ge(1-\eta)a_\infty-C_\eta$. Therefore, at fixed mass $M$,
\begin{equation}\label{eq:r10-coercivity}
 E_k(f)\ge\frac12\left[1-\eta-(M/M_\infty)^r\right]A_\infty(f)
             -\frac12C_\eta M.
\end{equation}
If $M<M_\infty$, choose $\eta$ to make the bracket positive. This
proves a finite lower bound and boundedness of minimizing sequences.
If also $M>M_0$, sharpness and amplitude scaling give a real smooth
mass-$M$ trial $f$ with $E_0(f)<0$. Its low-frequency dilation satisfies
\[
 E_k(f_\lambda)=\lambda^m E_0(f)+O(\lambda^{m+q})<0
\]
for sufficiently small $\lambda$, proving $I(M)<0$.

We record the compactness argument. Amplitude scaling gives
\[
 E_k(\sqrt\vartheta f)
 =\vartheta^{p/2}E_k(f)
   +\tfrac12(\vartheta-\vartheta^{p/2})[A_0(f)+B(f)]
 \qquad(\vartheta>1).
\]
Thus $I(\vartheta L)\le\vartheta I(L)$, with
$I(\vartheta L)\le\vartheta^{p/2}I(L)<\vartheta I(L)$ if $I(L)<0$,
whenever the masses are below $M_\infty$. For a mass $M$ with
$I(M)<0$, it follows that $I(M)/M<I(L)/L$ for every $0<L<M$.
Consequently
\begin{equation}\label{eq:r10-binding}
 I(M)<I(L)+I(M-L)\qquad(0<L<M).
\end{equation}
A bounded minimizing sequence cannot vanish locally in mass: the
localized GN inequality would then imply $\|f_n\|_p\to0$, contradicting
$I(M)<0$. After translations there is a nonzero weak $H^s$ limit $Q$
and almost-everywhere convergence. The quadratic forms split by weak
Hilbert convergence, and the nonlinear integral by the Brezis--Lieb
lemma. If $0<M(Q)<M$, normalize the remainder to its limiting mass;
boundedness of its kinetic and nonlinear terms makes the energy change
$o(1)$. Splitting then contradicts \eqref{eq:r10-binding}. Hence
$M(Q)=M$. For the remainder $w_n=f_n-Q$, the principal GN inequality
gives
\[
 \|w_n\|_p^p\le C M(w_n)^r A_0(w_n)\longrightarrow0.
\]
Energy splitting now forces $A_0(w_n)+B(w_n)\to0$, giving strong $H^s$
convergence and attainment.

Realness does not require a modulus inequality for a higher-order
operator. Write $f=X+iY$, put $\vartheta=M(X)/M(f)$, and, when both
components are nonzero, set $g=X/\sqrt\vartheta$,
$h=Y/\sqrt{1-\vartheta}$. These real functions have the same mass as
$f$. Evenness of $k$ splits its form over real and imaginary parts,
while convexity gives
\[
 |f|^p\le\vartheta|g|^p+(1-\vartheta)|h|^p,
 \qquad E_k(f)\ge\vartheta E_k(g)+(1-\vartheta)E_k(h).
\]
Therefore a real minimizing sequence exists and its strong limit is real.

Finally, for $M>M_\infty$ choose a smooth mass-$M$ function with
$E_\infty(f)<0$. High-frequency dilation and
\eqref{eq:r10-high-limit} give
\[
 \lambda^{-m}E_k(f_\lambda)\longrightarrow E_\infty(f)<0
 \qquad(\lambda\longrightarrow\infty),
\]
using $k\le C|\xi|^m$ for domination. Hence $I(M)=-\infty$.
\end{proof}

For each minimizer $U$ in the open window, the multiplier equation and
its scalar identity are
\begin{equation}\label{eq:r10-multiplier}
 k(D)U+\omega U-|U|^{2r}U=0,
 \qquad \omega M=r[A_0(U)+B(U)]-pI(M)>0.
\end{equation}
If additionally $k$ is $C^{N+2}$ off zero and satisfies the individual
order-$m$ derivative estimates
\[
 |\partial^\beta k(\xi)|\le C_\beta|\xi|^{m-|\beta|},
 \qquad |\beta|\le N+2,
\]
the elliptic bootstrap gives $U\in L^\infty\cap H^{m+1}$. A real
minimizer has zero momentum and its standing-wave formula gives a
global solution directly.

\subsection{The complete isotropic upper-endpoint classification}

Consider the interpolation
\begin{equation}\label{eq:r10-interpolation}
 h=a_\infty-a_0,\qquad
 k_q(\xi)=a_0(\xi)+h(\xi)\frac{|\xi|^q}{1+|\xi|^q}
         =a_\infty(\xi)-d_q(\xi),\qquad
 d_q(\xi)=\frac{h(\xi)}{1+|\xi|^q}.
\end{equation}
Write $D_q(f)=\int d_q|\widehat f|^2$. At mass $M_\infty$,
\begin{equation}\label{eq:r10-endpoint-energy}
 E_{k_q}(f)=E_\infty(f)-\tfrac12D_q(f),\qquad E_\infty(f)\ge0.
\end{equation}

\begin{theorem}[Upper-endpoint trichotomy]
\label{thm:r10-endpoint}
For \eqref{eq:r10-interpolation}, $q<m$ implies
$I(M_\infty)=-\infty$, whereas $q>m$ implies
$-\infty<I(M_\infty)<0$, attainment by a real minimizer, and
precompactness of every minimizing sequence modulo translations.
Both assertions allow anisotropic $a_0,a_\infty$.

If $a_0(\xi)=|\xi|^m$ and
$a_\infty(\xi)=(1+\varepsilon)|\xi|^m$, then at $q=m$,
\begin{equation}\label{eq:r10-borderline-endpoint}
 I(M_\infty)=-\tfrac12\varepsilon M_\infty,
 \qquad M_\infty=(1+\varepsilon)^{N/m}M_0,
\end{equation}
and the infimum is not attained.
\end{theorem}

\begin{proof}
Suppose $q<m$ and take an $a_\infty$ optimizer $Q$ of mass
$M_\infty$, so $E_\infty(Q)=0$. Then
\[
 \lambda^{-(m-q)}D_q(Q_\lambda)
 =\int\frac{h(\xi)}{\lambda^{-q}+|\xi|^q}|\widehat Q(\xi)|^2\,d\xi
 \longrightarrow
 \int h(\xi)|\xi|^{-q}|\widehat Q(\xi)|^2\,d\xi>0.
\]
The limit is finite: $0<m-q<m$, and the integrands are dominated by
$C|\xi|^{m-q}|\widehat Q|^2$. Thus
\eqref{eq:r10-endpoint-energy} tends to $-\infty$ along this sequence.

Suppose $q>m$. Then $d_q$ is bounded and tends to zero at infinity, so
\eqref{eq:r10-endpoint-energy} supplies the finite lower bound
$-\|d_q\|_\infty M_\infty/2$. Negativity follows from the
low-frequency trial since $M_\infty>M_0$. We must prove boundedness of
minimizing sequences because \eqref{eq:r10-coercivity} degenerates at
this endpoint. If instead a subsequence satisfies
$A_\infty(f_n)=t_n^m\to\infty$, write
$f_n(x)=t_n^{N/2}V_n(t_nx)$. Then
\[
 M(V_n)=M_\infty,\quad A_\infty(V_n)=1,\quad
 E_\infty(V_n)=t_n^{-m}[E_{k_q}(f_n)+D_q(f_n)/2]\longrightarrow0.
\]
Normalized sharp-GN compactness gives, after translations,
$V_n\to Q$ strongly in $H^s$, where $Q$ is an $a_\infty$ optimizer.
The bounded multiplier and strong $L^2$ convergence permit replacement
of $V_n$ by $Q$ in
\[
 D_q(f_n)=\int d_q(t_n\xi)|\widehat V_n(\xi)|^2\,d\xi.
\]
The expression for $Q$ tends to zero by dominated convergence. Thus
$D_q(f_n)\to0$ and \eqref{eq:r10-endpoint-energy} implies
$\liminf E_{k_q}(f_n)\ge0$, contradicting the negative infimum.
Minimizing sequences are therefore bounded. Amplitude scaling proves
\eqref{eq:r10-binding} also at $M=M_\infty$, since both positive
submasses lie strictly below it. The compactness and realness arguments
of Theorem~\ref{thm:r10-two-thresholds} now apply unchanged.

In the isotropic case $q=m$,
$d_m(\xi)=\varepsilon|\xi|^m/(1+|\xi|^m)<\varepsilon$.
Every admissible nonzero function consequently has
$E_{k_m}(f)>-\varepsilon M_\infty/2$. Conversely, high-frequency
dilations of an $a_\infty$ optimizer have zero homogeneous energy and
$D_m(Q_\lambda)\to\varepsilon M_\infty$ by dominated convergence.
This proves the exact infimum and its nonattainment.
\end{proof}

For completeness, the anisotropic borderline $q=m$ cannot be replaced
by the isotropic formula. Let $\mathcal G_\infty$ be the equality set
normalized by $M=M_\infty$, $A_\infty=1$. Define
\[
 J_{\rm esc}=-\frac12\sup_{Q\in\mathcal G_\infty}
       \int h(\xi/|\xi|)|\widehat Q(\xi)|^2\,d\xi.
\]
The supremum is attained by normalized GN compactness and continuity of
this bounded multiplier form. High-frequency optimizer trials give
$I(M_\infty)\le J_{\rm esc}$. If this inequality is strict,
minimizing sequences are bounded and attainment follows. Indeed an
unbounded sequence, normalized exactly as in the preceding proof,
has a limiting homogeneous optimizer and energy liminf at least
$J_{\rm esc}$. We make no assertion about attainment in the general
anisotropic equality case.

Every minimizer of \eqref{eq:r10-interpolation}, including the attained
upper endpoint $q>m$, obeys the exact identity
\begin{equation}\label{eq:r10-exact-negative-energy}
 E_{k_q}(U)=-\frac{q}{2m}
 \int h(\xi)\frac{|\xi|^q}{(1+|\xi|^q)^2}|\widehat U(\xi)|^2\,d\xi<0.
\end{equation}
To prove it, differentiate the explicit Fourier energy of the
mass-preserving dilation at its scalar minimum $\lambda=1$. The
derivative is dominated by $C|\xi|^m$. Subtract the resulting identity
divided by $m$ from the energy. This uses no spatial moment.

\subsection{Sharp behavior at the lower threshold}

Return to the general symbol and assume additionally, uniformly on
$S^{N-1}$,
\begin{equation}\label{eq:r10-low-expansion}
 R^{-(m+q)}b(R\theta)\longrightarrow c_*(\theta)
 \qquad(R\downarrow0),
\end{equation}
where $c_*$ is positive elliptic and homogeneous of degree $m+q$.
Let
\[
 \mathcal G_0=\{Q:M(Q)=M_0,\ A_0(Q)=1,\ \|Q\|_p^p=p/2\},
 \qquad C_*(f)=\int c_*(\xi)|\widehat f(\xi)|^2\,d\xi.
\]
Assume that at least one $Q\in\mathcal G_0$ has $C_*(Q)<\infty$.
This finite-moment assumption is automatic for $0<q\le m+2$ under the
stated regularity of $a_0$: its optimizer equation and the elliptic
bootstrap give $Q\in H^{m+1}$. For larger $q$ it is an explicit
hypothesis; no unlimited smoothness of a non-polynomial nonlinearity is
presumed.

Define
\begin{equation}\label{eq:r10-selection-constants}
 \begin{gathered}
 C_{\min}=\min_{Q\in\mathcal G_0}C_*(Q),\qquad
 \beta_* =\left[\frac{m}{(m+q)C_{\min}}\right]^{1/q},\qquad
 K_* =\frac{q}{2(m+q)}\beta_*^m,\\
 \mathcal G_{\rm sel}=\{Q\in\mathcal G_0:C_*(Q)=C_{\min}\}.
 \end{gathered}
\end{equation}
These definitions are well posed. The assumed finite trial gives
finiteness; positivity follows from Jensen's inequality,
\[
 \int|\xi|^{m+q}|\widehat Q|^2
 \ge M_0^{-q/m}\left(\int|\xi|^m|\widehat Q|^2\right)^{1+q/m},
\]
and $A_0(Q)=1$. A minimizing sequence in $\mathcal G_0$ is compact in
$H^s$ modulo translations, and the nonnegative $C_*$ form is lower
semicontinuous. Hence the minimum is attained. The selected set is
compact in $H^s$ modulo translations and constant phase.

\begin{theorem}[Scale law and optimizer selection]
\label{thm:r10-selection}
Let $M\downarrow M_0$ and let $U_M$ be any minimizer of $I(M)$. Set
\begin{equation}\label{eq:r10-normalization}
 \alpha=\frac{M}{M_0},\qquad d=\alpha^r-1,\qquad
 \lambda_M=\left[\frac{A_0(U_M)}{\alpha}\right]^{1/m},\qquad
 U_M(x)=\sqrt\alpha\,\lambda_M^{N/2}V_M(\lambda_Mx).
\end{equation}
Then $M(V_M)=M_0$, $A_0(V_M)=1$, and
\begin{align}
 I(M)&=-\alpha K_*d^{1+m/q}(1+o(1)),\label{eq:r10-energy-law}\\
 \lambda_M/d^{1/q}&\longrightarrow\beta_*,\label{eq:r10-scale-law}\\
 \operatorname{dist}_{H^s}(V_M,\mathcal G_{\rm sel})&\longrightarrow0,
       \label{eq:r10-shape-law}\\
 1-\frac2p\|V_M\|_p^p&=o(d),\label{eq:r10-deficit-law}\\
 \frac{B(U_M)}{A_0(U_M)}&=\frac{m}{m+q}d+o(d).
       \label{eq:r10-equipartition-law}
\end{align}
The distance includes translations and constant phases, already
contained in $\mathcal G_{\rm sel}$. For the multiplier in
\eqref{eq:r10-multiplier},
\begin{equation}\label{eq:r10-frequency-law}
 \omega_M=\frac{r}{M_0}\beta_*^m d^{m/q}(1+o(1)).
\end{equation}
\end{theorem}

\begin{proof}
For the normalization $M(V)=M_0$, $A_0(V)=1$, let
$H(V)=1-(2/p)\|V\|_p^p\ge0$. The exact energy formula is
\begin{equation}\label{eq:r10-exact-scaled-energy}
 I(M)=\frac{\alpha\lambda_M^m}{2}
          [-d+(1+d)H(V_M)]
       +\frac\alpha2\int b(\lambda_M\xi)|\widehat V_M(\xi)|^2\,d\xi.
\end{equation}
Negativity implies $0\le H(V_M)<d/(1+d)$. Therefore normalized
sharp-GN compactness applies, giving subsequential convergence, modulo
translation and phase, to some $Q\in\mathcal G_0$ strongly in $H^s$.

First $\lambda_M\to0$. Otherwise along a subsequence
$\lambda_M\ge\ell>0$, and \eqref{eq:r10-bounds} gives
\[
 \frac{B(U_M)}{A_0(U_M)}
 \ge c\int\min\{\lambda_M^q|\xi|^{m+q},|\xi|^m\}
                         |\widehat V_M(\xi)|^2\,d\xi.
\]
Choose a compact annulus carrying positive Fourier mass of the limiting
$Q$. The right side is then bounded below by a fixed positive number.
But \eqref{eq:r10-exact-scaled-energy} implies $B/A_0<d\to0$, a
contradiction.

Define the rescaled correction form
\[
 D_\lambda(V)=\int\frac{b(\lambda\xi)}{\lambda^{m+q}}
                              |\widehat V(\xi)|^2\,d\xi.
\]
Whenever $\lambda_n\downarrow0$ and $V_n\to Q$ strongly in $H^s$,
\begin{equation}\label{eq:r10-correction-liminf}
 \liminf_n D_{\lambda_n}(V_n)\ge C_*(Q).
\end{equation}
Indeed on a compact annulus, \eqref{eq:r10-low-expansion} gives uniform
multiplier convergence and strong Fourier $L^2$ convergence gives
convergence of the integral. Discard the complement by nonnegativity,
then exhaust all frequencies. In particular the liminf is at least
$C_{\min}$ for the normalized minimizing profiles. This step does not
assume uniform higher moments of those profiles.

For a fixed $Q\in\mathcal G_{\rm sel}$,
\eqref{eq:r10-bounds} bounds the rescaled multiplier by
$C|\xi|^{m+q}$. Dominated convergence gives
$D_\lambda(Q)\to C_{\min}$. Testing with
$\sqrt\alpha(\beta_*d^{1/q})^{N/2}Q(\beta_*d^{1/q}x)$ proves
\begin{equation}\label{eq:r10-energy-limsup}
 \limsup_{M\downarrow M_0}\frac{I(M)}{\alpha d^{1+m/q}}
 \le-K_*<0.
\end{equation}
On the other hand, negativity in
\eqref{eq:r10-exact-scaled-energy} gives
$\lambda_M^qD_{\lambda_M}(V_M)<d$. Equation
\eqref{eq:r10-correction-liminf} consequently yields
$\lambda_M^q\le2d/C_{\min}$ for small $d$. The lower bound
$I(M)\ge-\alpha d\lambda_M^m/2$, combined with
\eqref{eq:r10-energy-limsup}, yields
$\lambda_M^m\ge K_*d^{m/q}$ for small $d$. Thus
\begin{equation}\label{eq:r10-scale-comparable}
 0<c\le\beta_M:=\lambda_M/d^{1/q}\le C<\infty.
\end{equation}
The same inequalities bound $D_{\lambda_M}(V_M)$ and $H(V_M)/d$.

Extract a subsequence with $\beta_M\to\beta\in(0,\infty)$,
$V_M\to Q\in\mathcal G_0$ strongly in $H^s$ modulo symmetries, and
convergent bounded scalar quantities $D_{\lambda_M}(V_M)$ and
$H(V_M)/d$. Divide \eqref{eq:r10-exact-scaled-energy} by
$\alpha d^{1+m/q}$:
\begin{equation}\label{eq:r10-final-scaled-energy}
 \frac{I(M)}{\alpha d^{1+m/q}}
 =\frac12\left[-\beta_M^m
 +(1+d)\beta_M^m\frac{H(V_M)}d
 +\beta_M^{m+q}D_{\lambda_M}(V_M)\right].
\end{equation}
Its liminf is at least
\[
 \tfrac12[-\beta^m+C_*(Q)\beta^{m+q}]
 \ge\min_{t>0}\tfrac12[-t^m+C_{\min}t^{m+q}]=-K_*.
\]
Equality with \eqref{eq:r10-energy-limsup} forces
$\beta=\beta_*$, the unique scalar minimizer, and
$C_*(Q)=C_{\min}$. It also forces
$H(V_M)/d\to0$ and $D_{\lambda_M}(V_M)\to C_{\min}$.
Every extracted limit has these properties. This proves
\eqref{eq:r10-energy-law}--\eqref{eq:r10-deficit-law}, including the
distance statement. Finally,
\[
 \frac{B(U_M)}{A_0(U_M)}
 =\lambda_M^qD_{\lambda_M}(V_M)
 =d\beta_*^qC_{\min}+o(d)=\frac{m}{m+q}d+o(d).
\]
The multiplier identity \eqref{eq:r10-multiplier}, together with
$A_0(U_M)=\alpha\lambda_M^m$, $B/A_0=O(d)$, and $I/A_0=O(d)$, gives
$\omega_M/\lambda_M^m\to r/M_0$, proving
\eqref{eq:r10-frequency-law}.
\end{proof}

The proof also gives the normalized frequency correction
\[
 \frac{\omega_M}{\lambda_M^m}
 =\frac r{M_0}+\frac1{M_0}\left(r+\frac q{m+q}\right)d+o(d).
\]
This does not assert an additional-order expansion of $\lambda_M$.

\subsection{The limiting functional}

The selection has a precise variational interpretation. For $d>0$ set
$\alpha_d=(1+d)^{1/r}$ and $\ell=d^{1/q}$, and define, on $H^s$,
\[
 \mathcal F_d(W)=
 \frac{E_k(\sqrt{\alpha_d}\,\ell^{N/2}W(\ell\,\cdot))}
      {\alpha_d d^{1+m/q}}
 \quad\hbox{if }M(W)=M_0,
 \qquad \mathcal F_d(W)=+\infty\quad\hbox{otherwise}.
\]
The exact expression on the mass constraint is
\begin{equation}\label{eq:r10-gamma-exact}
 \mathcal F_d(W)=\frac{E_0(W)}d-\frac1p\|W\|_p^p
       +\frac12\int\frac{b(\ell\xi)}{\ell^{m+q}}
                                 |\widehat W(\xi)|^2\,d\xi.
\end{equation}

\begin{proposition}[Variational selection on the equality manifold]
\label{prop:r10-gamma}
In the strong $H^s$ topology, $\mathcal F_d$ $\Gamma$-converges as
$d\downarrow0$ to
\[
 \mathcal F_0(W)=
 \begin{cases}
 \tfrac12[C_*(W)-A_0(W)],&M(W)=M_0,\ E_0(W)=0,\\
 +\infty,&\text{otherwise}.
 \end{cases}
\]
Its minimizers are exactly the dilates
$\beta_*^{N/2}Q(\beta_*\,\cdot)$ with
$Q\in\mathcal G_{\rm sel}$, including their symmetries.
\end{proposition}

\begin{proof}
On the fixed mass constraint $E_0\ge0$. If $W_n\to W$ strongly in
$H^s$ and $E_0(W)>0$, continuity of the homogeneous energy and the
nonlinear norm makes \eqref{eq:r10-gamma-exact} tend to $+\infty$.
If $E_0(W)=0$, discard its nonnegative first term and apply the annular
liminf argument \eqref{eq:r10-correction-liminf}. Since equality in
the GN inequality gives $\|W\|_p^p/p=A_0(W)/2$, this is the required
liminf. A limit of the fixed mass constraint has mass $M_0$ by strong
$L^2$ convergence. For finite $\mathcal F_0(W)$, the constant sequence
$W$ is a recovery sequence, using domination by
$C|\xi|^{m+q}|\widehat W|^2$. This proves the $\Gamma$-limit.
Every nonzero equality profile is a dilation of an element of
$\mathcal G_0$; minimizing its scalar dilation gives exactly
\eqref{eq:r10-selection-constants}.
\end{proof}

Translation invariance prevents global equicoercivity in $H^s$; none
is asserted. Theorem~\ref{thm:r10-selection} separately proves
compactness of the rescaled minimizers modulo translations. The
limiting principle says that the first nonzero low-frequency correction
selects a subset of the entire homogeneous optimizer manifold.

\subsection{A joint small-perturbation and mass-excess limit}

There is a complementary selection limit that needs no higher-moment
hypothesis. Fix a positive elliptic homogeneous symbol $a$ as above and
let
\[
 b_0(\xi)=a(\xi)\frac{|\xi|^q}{1+|\xi|^q},\qquad
 k_\varepsilon=a+\varepsilon b_0,\qquad M_c=M_{c,a}.
\]
Write $A=A_a$, $B_0(f)=\int b_0|\widehat f|^2$, and let
$I_\varepsilon$ be the constrained infimum for $k_\varepsilon$. The
entire homogeneous equality manifold, with no kinetic normalization, is
\[
 \mathcal H=\{Q:M(Q)=M_c,\ E_a(Q)=0\}.
\]
For $0<\tau<1$ define
\[
 J(\tau)=\inf_{Q\in\mathcal H}\tfrac12[B_0(Q)-\tau A(Q)],\qquad
 \mathcal S_\tau=\operatorname*{argmin}_{Q\in\mathcal H}
                                \tfrac12[B_0(Q)-\tau A(Q)].
\]

\begin{proposition}[Selection by an arbitrarily weak perturbation]
\label{prop:r10-crossover}
For every $0<\tau<1$, $J(\tau)$ is finite, negative, and attained, and
$\mathcal S_\tau$ is compact in $H^s$ modulo translations and phase.
Let $\varepsilon\downarrow0$, set
\[
 \alpha_\varepsilon=(1+\varepsilon\tau)^{1/r},\qquad
 M_\varepsilon=\alpha_\varepsilon M_c,
\]
and let $U_\varepsilon$ be any minimizer of
$I_\varepsilon(M_\varepsilon)$. Then
\[
 \frac{I_\varepsilon(M_\varepsilon)}
             {\alpha_\varepsilon\varepsilon}\longrightarrow J(\tau),
 \qquad
 \operatorname{dist}_{H^s}
       (U_\varepsilon/\sqrt{\alpha_\varepsilon},\mathcal S_\tau)
       \longrightarrow0.
\]
No spatial rescaling is used. The same assertions hold along
$\tau_\varepsilon\to\tau\in(0,1)$, with
$\alpha_\varepsilon=(1+\varepsilon\tau_\varepsilon)^{1/r}$ and limiting
set $\mathcal S_\tau$; the compactness estimates are uniform for
$\tau_\varepsilon$ in a compact subinterval of $(0,1)$.
\end{proposition}

\begin{proof}
For every $\eta>0$,
$B_0(f)\ge(1-\eta)A(f)-C_\eta M(f)$. Thus the objective defining
$J(\tau)$ is coercive in $A$ on $\mathcal H$ if $\eta<1-\tau$.
Low-frequency dilation of any homogeneous optimizer makes the
objective negative. Its minimizing sequences consequently have $A$
bounded above and bounded away from zero, the latter because the
objective is at least $-\tau A/2$. Normalize $A$ to one and apply
sharp-GN compactness. The normalizing scales stay in a compact subset
of $(0,\infty)$, giving strong $H^s$ convergence modulo translations.
Continuity of the order-$m$ form $B_0$ proves attainment and compactness
of the selected set.

On the mass constraint $M(f)=M_c$, the exact identity is
\begin{equation}\label{eq:r10-crossover-functional}
 \frac{E_{k_\varepsilon}(\sqrt{\alpha_\varepsilon}f)}
             {\alpha_\varepsilon\varepsilon}
 =\frac{E_a(f)}\varepsilon+\frac12B_0(f)
                              -\frac\tau p\|f\|_p^p.
\end{equation}
It $\Gamma$-converges in the strong $H^s$ topology to
$[B_0-\tau A]/2$ on $\mathcal H$ and $+\infty$ off $\mathcal H$:
use $E_a\ge0$, continuity of the last two terms, and constant equality
profiles as recovery sequences.

We prove compactness of actual minimizers separately. The sharp GN
inequality at mass $\alpha_\varepsilon M_c$ gives
\[
 E_{k_\varepsilon}(U)
 \ge\frac\varepsilon2
       [(1-\eta-\tau)A(U)-C_\eta M(U)].
\]
Testing \eqref{eq:r10-crossover-functional} with a selected profile
gives the exact bound
$I_\varepsilon/(\alpha_\varepsilon\varepsilon)\le J(\tau)<0$.
Hence $A(U_\varepsilon)$ is uniformly bounded. The other lower bound
$E_{k_\varepsilon}(U)\ge-\varepsilon\tau A(U)/2$ forces it to be
bounded away from zero. With
$f_\varepsilon=U_\varepsilon/\sqrt{\alpha_\varepsilon}$, identity
\eqref{eq:r10-crossover-functional} now gives
$E_a(f_\varepsilon)=O(\varepsilon)$. Normalize the kinetic form and
use sharp-GN compactness to obtain a strong $H^s$ limit
$Q\in\mathcal H$ modulo translation, without any escaping spatial
scale. The $\Gamma$-liminf is at least $J(\tau)$ and the trial upper
bound is at most $J(\tau)$, proving convergence of the energies and
selection of $Q\in\mathcal S_\tau$.

For $\tau_\varepsilon\to\tau$, the same proof has uniform constants:
choose $\eta$ below the distance of the compact parameter interval to
one, and a single sufficiently low-frequency optimizer trial that is
strictly negative for the smallest parameter in that interval.
Continuity of the bounded-energy objective gives the stated limit.
\end{proof}

Along every convergent subsequence in this proposition, the multiplier
satisfies $\omega_\varepsilon\to rA(Q)/M_c$ for
$Q\in\mathcal S_\tau$. Thus a perturbation tending to zero can select
a finite spatial scale and frequency from the critical homogeneous
family when the mass excess has the same order as the perturbation.
The selection functional $J$ is concave and strictly decreasing. The
first property follows by taking an infimum of affine functions; for
the second, test the objective at $\tau_2>\tau_1$ with a minimizer
$Q_1$ at $\tau_1$ to obtain
\[
 J(\tau_2)\le J(\tau_1)-\tfrac12(\tau_2-\tau_1)A(Q_1)<J(\tau_1).
\]

\subsection{A quartic example with explicit exponents}

For $a_0(\xi)=|\xi|^4$ and
$b(\xi)=\varepsilon|\xi|^6/(1+|\xi|^2)$, let
\[
 S_6=\min_{Q\in\mathcal G_0}\int|\xi|^6|\widehat Q(\xi)|^2\,d\xi,
 \qquad d=(M/M_0)^{4/N}-1,\qquad \alpha=M/M_0.
\]
The higher-moment hypothesis holds, and $C_{\min}=\varepsilon S_6$.
Theorem~\ref{thm:r10-selection} gives
\[
 I(M)\sim-\alpha\frac{2}{27\varepsilon^2 S_6^2}d^3,
 \qquad \lambda_M^2\sim\frac{2}{3\varepsilon S_6}d,
 \qquad
 \omega_M\sim\frac{4}{NM_0}\frac{4}{9\varepsilon^2 S_6^2}d^2.
\]
Thus the spatial dilation length grows as $d^{-1/2}$, the principal
kinetic energy vanishes as $d^2$, and the negative constrained energy
vanishes as $d^3$. The selected critical profiles minimize their sixth
Fourier moment. At the other end, $q=2<m=4$ implies
$I(M_\infty)=-\infty$.

For the isotropic family as a whole the result is a complete constrained
energy classification: zero nonattained infimum up to $M_0$, negative
attained infimum in $(M_0,M_\infty)$, and $-\infty$ above
$M_\infty$, with the upper endpoint determined by the three cases in
Theorem~\ref{thm:r10-endpoint}. The new asymptotic concerns a class of
nonhomogeneous equations; it does not resolve the unrestricted dynamics
of the original homogeneous BNLS.
